% Plantilla latex para protocolo de tesis de posgrado en la Facultad de Ingeniería, Universidad Autónoma de Querétaro.
%
% @author {Edgar Lara Arellano}
% @email {elara28@alumnos.uaq.mx}
% @year 2024
% @version 1.0
% @license CC-BY-NC
%
% Se agradecen sus comentarios, contribuciones,  reporte de bugs y la difusión de esta plantilla


% Logos tomados de: https://ingenieria.uaq.mx/index.php/descargas
% Elementos complementarios tomados de: https://www.uaq.mx/f_ingenieria/images/Assets_web_ing/PAQUETE_GRAFICO_FI_PARA_DESCARGA/Manual_identidad_oficil_fi/MANUAL-OFICIAL-FI-2019.pdf

%Plantilla basada en los requerimientos planteados en [Guía para la escritura de Tésis, FI, UAQ](https://dip.uaq.mx/docs/posgrado/Guia-tesis-CNyE.pdf)



\documentclass[12pt, letterpaper, spanish, twoside]{article}
\AddToHook{cmd/section/before}{\newpage}
\usepackage{./common/UAQ}

\graphicspath{{./figures/}}


\begin{document}


\input{./chapters/portada_externa}


\backgroundsetup{
	scale = 1,
	angle = 0,
	opacity = 1,
	contents = {
		\includegraphics[width=\paperwidth]{figures/externalCoverBackground.pdf}
	}
}
\input{./chapters/portada_interna}




\input{./chapters/0_dedicatorias}

\newpage
\input{./chapters/0_agradecimientos}

\newpage
\input{./chapters/0_abreviaturas}

\newpage
\input{./chapters/0_resumen}
\newpage



\pagenumbering{\Roman}
\setlength{\cftsecnumwidth}{3em} %Se cambia el espacio entre el item y el nombre de la seccion
\setlength{\cftsubsecindent}{1.3cm}
\setlength{\cftsubsubsecindent}{2.4cm}
\tableofcontents
\listoftables
\listoffigures




\pagestyle{plain}
\pagenumbering{arabic}




\section{Introducción/planteamiento del problema y justificación}
\input{./chapters/1_introduccion}

\section{Antecedentes}
\input{./chapters/2_antecedentes}

\section{Hipótesis}
\input{./chapters/3_hipotesis}

\section{Objetivos}
\input{./chapters/4_objetivos}

\section{Material y Métodos o Metodología}
\input{./chapters/5_metodologia}

\section{Resultados y discusión}
\input{./chapters/6_resultados}

\section{Conclusiones}
\input{./chapters/7_conclusiones}



\nocite{*} % Para mostrar todas las refencias sin citarlas
\renewcommand{\refname}{Referencias bibliográficas}
\bibliographystyle{IEEEtran}
%\bibliographystyle{apalike}
\bibliography{referencias}

\end{document}